\documentclass[
]{ajs}

%% recommended packages
\usepackage{orcidlink,thumbpdf,lmodern}

\usepackage[utf8]{inputenc}

\author{
FirstName LastName\\University/Company \And Second Author\\Affiliation
\AND Third Author\\Universitat Autònoma\\
de Barcelona
}
\title{A Capitalized Title: Something about a Topic}

\Plainauthor{FirstName LastName, Second Author, Third Author}
\Plaintitle{A Capitalized Title: Something about a Topic}
\Shorttitle{A Capitalized Title}


\Abstract{
This guide of the Austrian Journal of Statistics (AJS) shows the
required style of manuscripts submitted to the Austrian Journal of
Statistics. We strongly recommend to follow the given instructions
exactly already when writing the paper. When submitting papers to AJS,
it is essential to follow these guidelines. Manuscripts submitted
without consideration of the AJS style guide will most likely not be
forwarded for review, but rejected with the possibility of resubmission.
}

\Keywords{keywords, not capitalized, \proglang{Java}}
\Plainkeywords{keywords, not capitalized, Java}

%% publication information
%% \Volume{50}
%% \Issue{9}
%% \Month{June}
%% \Year{2012}
%% \Submitdate{}
%% \Acceptdate{2012-06-04}

\Address{
    FirstName LastName\\
    University/Company\\
    First line\\
Second line\\
  E-mail: \email{name@company.com}\\
  URL: \url{https://posit.co}\\~\\
        Third Author\\
    Universitat Autònoma de Barcelona\\
    Department of Statistics and Mathematics,\\
Faculty of Biosciences,\\
Universitat Autònoma de Barcelona\\
  
  
  }


% tightlist command for lists without linebreak
\providecommand{\tightlist}{%
  \setlength{\itemsep}{0pt}\setlength{\parskip}{0pt}}



\usepackage{amsmath}



\begin{document}



\section{Placeholder!! can be deleted}\label{placeholder-can-be-deleted}

\section{Data and Methods}\label{data-and-methods}

\subsection{Data Description}\label{data-description}

This project is based on genetic and osteological data obtained from
human remains recovered from Mycenaean collective burial contexts in
Greece, dated to approximately 1500--1000 BCE. The material derives from
two archaeological sites: Amfissa, which consists of a single tholos
tomb, and Elateia, a large chamber tomb complex containing 85 chambers.
From Elateia, eight chambers were selected for detailed genetic
investigation. DNA preserved in skeletal material was extracted and
analyzed to infer biological relationships among individuals buried
within and between tombs.

The analysis draws on three interconnected datasets that capture
information at different levels of resolution. One dataset summarizes
characteristics at the tomb level. It includes variables such as the
minimum number of individuals identified, the number of individuals
successfully analyzed genetically, the estimated duration over which the
tomb was used, the proportion of successful DNA samples, the total
number of samples taken, and the distribution of sampled skeletal
elements (including petrous bones, teeth, talus or phalanx, and other
elements). These variables highlight differences between tombs in
preservation, sampling, and burial practices.

A second dataset contains metadata describing the sampled individuals.
Each entry represents a single individual and records the associated
tomb, estimated biological sex, age category, and the skeletal element
used for DNA extraction. This information provides the biological and
contextual background necessary to evaluate how individual
characteristics relate to genetic data quality and inferred
relationships.

The third dataset reports pairwise genetic relationships between
individuals based on ancient DNA analyses. For each pair, it includes
estimates of genetic relatedness, the number of overlapping single
nucleotide polymorphisms (SNPs), and a categorical classification of
relationship degree following established criteria. Since individuals
can be involved in multiple pairwise comparisons, observations in this
dataset are not independent, resulting in a complex relational
structure.

\subsection{Data Preparation for
Model}\label{data-preparation-for-model}

Prior to statistical modeling, the raw datasets were processed and
integrated to obtain consistent pairwise datasets suitable for analysis.
Data preparation involved filtering relevant observations, harmonizing
variable definitions, combining information across different levels, and
constructing derived variables required for modeling.

At the individual level, metadata were first restricted to the burial
sites and tombs included in the analysis. Individuals from Lapoutsi were
excluded, and only samples from Amfissa tholos and selected tombs at
Elateia were retained. Age estimates were harmonized by collapsing
ambiguous or overlapping categories into two broad groups, namely
subadult and adult or undefined, to ensure sufficient sample sizes per
category.

Tomb-level parameters were prepared separately and standardized to
ensure consistency across datasets. These parameters include the minimum
number of individuals, the number of genetically analyzed individuals,
the estimated duration of tomb use, and the proportion of successful DNA
samples. To allow analyses at different spatial resolutions, an
additional set of tomb-level parameters was constructed in which Tombs
46, 56, and 62 at Elateia were treated as a single combined burial unit.
Corresponding values for this grouped tomb were derived from the
original data and added alongside the individual tomb records.

Pairwise genetic kinship data were then processed to generate the main
analytical datasets. Relationship classifications were cleaned by
distinguishing between confidently assigned relationships and uncertain
cases. Pairs with insufficient genetic information were marked as
uncertain, and a binary relatedness indicator was derived that
distinguishes related and unrelated pairs while excluding uncertain
observations. In addition, a global relationship classification based on
a stricter minimum SNP overlap threshold of 15,000 was created to reduce
bias associated with low-quality genetic data.

Individual-level metadata were merged into the kinship dataset for both
members of each pair, adding information on tomb affiliation, sex, age
category, and sampled skeletal element. Based on these merged data, two
parallel analytical datasets were constructed. The first dataset retains
the original tomb structure, in which Tombs 46, 56, and 62 are treated
as separate burial contexts. The second dataset includes an additional
set of observations in which pairs involving individuals from these
tombs are duplicated and reassigned to a combined tomb category. This
approach preserves information at the individual tomb level while
enabling sensitivity analyses that assess the impact of treating closely
related burial contexts as a single unit.

Finally, both datasets were augmented with tomb-level characteristics by
linking each individual's tomb to the corresponding sampling success,
number of analyzed individuals, and length of use. From these variables,
pairwise summaries were derived, including the minimum sampling success,
the minimum number of analyzed individuals, and the absolute difference
in tomb use duration between the two burial contexts. Additional
pair-level variables were constructed to describe sex composition, age
composition, shared tomb membership, combined skeletal element types,
and an indicator of bone quality based on the presence of petrous bone
samples.

The resulting datasets consist of pairwise observations enriched with
individual- and tomb-level information and derived covariates. Together,
they allow the evaluation of genetic relatedness patterns under
alternative representations of burial organization and provide the basis
for the multi-membership modeling strategy applied in the subsequent
analysis.

\subsection{The Multi-Membership Generalized Linear Mixed
Model}\label{the-multi-membership-generalized-linear-mixed-model}

\subsubsection{Model Definition and
Motivation}\label{model-definition-and-motivation}

To analyze the kinship patterns within the dataset, we require a
modeling framework that can account for the inherent dependencies within
our observations. In this study, an observation is defined as a pair of
individuals \((i,j)\), where the response variable \(Y_{ij}\) indicates
whether the pair shares a biological relationship.

\textbf{Motivation: The Complexity of Pairwise Data}

The data structure presents two primary challenges that traditional
regression models cannot handle: Multiple Membership of Individuals:
Each individual appears in multiple pairs. For example, if individual
\(i\) is compared to \(j\), \(k\), and \(l\), the errors associated with
those three observations will be correlated because they all share the
unique genetic and environmental traits of individual \(i\). Tomb-Level
Dependencies: Pairs can exist both within the same tomb or across
different tombs. Because individuals within a tomb may share common
environmental or social factors, we must account for the random ``tomb
effect.'' A pair \({i,j}\) effectively ``belongs'' to the tombs \(T(i)\)
and \(T(j)\).

To address this, we utilize a Multi-Membership Generalized Linear Mixed
Model (MM-GLMM). This model is specifically designed for cases where an
observation does not belong to a single cluster, but rather ``members''
multiple clusters (in our case, two individuals and two tombs)
simultaneously \citet{browne2001multiple}.

\textbf{Model Definition}

Let \(Y_{ij}\) be a binary indicator where \(Y_{ij} = 1\) if individuals
\(i\) and \(j\) are related, and \(Y_{t,ij} = 0\) otherwise. We assume
\(Y_{ij}\) follows a Bernoulli distribution:

\[
Y_{t,ij} \sim \text{Bernoulli}(\pi_{ij})
\]

where \(\pi_{ij} = P(\text{related}_{ij} = 1)\) is the probability of a
relationship existing between the pair. We model this probability using
a logit link function to accommodate the complex random effects
structure:

\[
\text{logit}(\pi_{ij}) = \alpha + \mathbf{X}_{ij}^{\top}\boldsymbol{\beta} 
    + \underbrace{(w_i u_i + w_j u_j)}_{\text{Individual RE}} 
    + \underbrace{(w_{T(i)} v_{T(i)} + w_{T(j)} v_{T(j)})}_{\text{Tomb RE}}
\]

Components of the Model:

\begin{itemize}
\tightlist
\item
  \(\alpha\): The global intercept, representing the baseline log-odds
  of two individuals being related.
\item
  \(\mathbf{X}_{ij} \beta\): The fixed effects component, where
  \(\mathbf{X}_{ij}\) contains pair-level covariates.
\item
  \(u_i, u_j\): Random effects for the two individuals, where
  \(u \sim \mathcal{N}(0, \sigma^2_{\text{ind}})\)
\item
  \(v_{T(i)}, v_{T(j)}\): Random effects for the tombs associated with
  individuals \(i\) and \(j\), where
  \(v \sim \mathcal{N}(0, \sigma^2_{\text{tomb}})\)
\item
  \(w\): Weights assigned to the random effects. Following standard
  multi-membership notation \citet{browne2001multiple}, we set equal
  weights \(w_i = w_j = 0.5\) and \(w_{T(i)} = w_{T(j)} = 0.5\). This
  ensures that the total variance contributed by the individual or tomb
  level remains consistent and equal.
\end{itemize}

\subsubsection{Fitting the Model with
brms}\label{fitting-the-model-with-brms}

Following the theoretical definition of the MM-GLMM, this section
details the practical implementation, predictor selection, and the
Bayesian framework used for estimation.

\emph{Predictor Selection and the Causal Framework} The selection of
covariates is guided by the Directed Acyclic Graph (DAG) presented in
Section {[}2.3.1. DAG{]}. To estimate the probability of relatedness, we
focus on variables situated along the social, biological branches of our
causal model. Specifically, the following predictors are included in the
model: Pairwise Sex Combination, Pairwise Age Category, and Same Tomb
Indicator.

Notably, DNA preservation variables (Number of Overlapping SNPs,
Analyzed skeletal element,\ldots) are excluded from the model. While
preservation significantly influences data quality and the initial
ability to identify SNPs, it does not exert a causal influence on true
biological relatedness itself. Because a rigorous quality filter, a
15,000 SNP threshold, was applied during data preprocessing, the
influence of preservation bias is effectively mitigated. Once a pair
meets this threshold, variation in kinship outcomes is no longer
explained by preservation quality.

\emph{Bayesian Implementation via brms} The model is fitted using the
brms package in R. The brms package acts as an interface to Stan,
utilizing Hamiltonian Monte Carlo (HMC) algorithms to perform Bayesian
inference. This approach is particularly robust for Multi-Membership
models, as it allows for flexible specification of complex random effect
structures that often struggle to estimate using frequentist Maximum
Likelihood methods \citet{burkner2017advanced}.

Two distinct versions of the model are implemented

Separated Tomb Model: Every tomb is treated as an individual level
within the random effect \(v_t\). Grouped Tomb Model: In this version,
three specific tombs, Elateia T46, T56, and T62, are treated as a single
pooled tomb.

\section{Result}\label{result}

Following the model specification, this chapter presents the findings
from the Bayesian Multi-Membership Generalized Linear Mixed Model.

\subsection{Model Diagnostics}\label{model-diagnostics}

Before interpreting the relationship between burial patterns and
kinship, it is essential to assess the ``sampling health'' of the model.
Because the MM-GLMM was fitted using Hamiltonian Monte Carlo (HMC) via
Stan, we evaluate the reliability of the posterior distributions using
two primary metrics: the potential scale reduction factor (\(\hat{R}\))
and the Effective Sample Size (ESS).

\textbf{Convergence and Chain Stability}: A primary indicator of model
health is the \(\hat{R}\) (R-hat) statistic, which compares the variance
within individual MCMC chains to the variance between chains. For all
parameters in our model, the \(\hat{R}\) values were consistently
\textbf{1.00}. Following modern Bayesian standards, an
\(\hat{R} < 1.01\) indicates that all four chains have successfully
converged to the same posterior distribution, confirming that the model
has reached stability and is not sensitive to the starting values of the
chains \citet{vehtari2021rank}.

\textbf{Sampling Efficiency}: The Effective Sample Size (ESS) provides
an estimate of the number of independent draws contained in the
posterior distribution, accounting for autocorrelation between
consecutive MCMC steps. In accordance with the standards implemented in
the brms framework, we report two measures:

\begin{itemize}
\tightlist
\item
  Bulk-ESS: Measures the sampling efficiency for the mean and median of
  the distribution.
\item
  Tail-ESS: Measures the efficiency for the tails of the distribution,
  which is critical for the reliability of the 95\% credible intervals.
\end{itemize}

As shown in the model output, both the Bulk-ESS and Tail-ESS for all key
predictors are well above the recommended thresholds (typically
\textgreater400 per chain). Many parameters reached between 3,000 and
22,000 independent samples. For instance, the ``Same Tomb'' predictor
achieved a Bulk-ESS of 22,386, indicating an exceptionally high degree
of precision in the estimation of the interment effect. These metrics
provide a high degree of confidence that the posterior estimates
discussed in the following sections are robust and represent truly
independent information from the data \citet{vehtari2021rank}.

\subsection{Model Output}\label{model-output}

Following the diagnostic verification, we present the posterior
estimates for the two versions of our model: the Separated Tomb Model (8
tombs) and the Grouped Tomb Model (6 tombs). Both versions show
remarkable stability in their regression coefficients, confirming that
the pooling of the Elateia tombs does not fundamentally alter the
estimated relationships. For a visual representation of the full effect
estimations and their associated 95\% Credible Intervals (CI), please
refer to the plot in ??Figure.

A primary observation is the dominance of Same Tomb Indicator, which
yielded the highest estimate in both models (2.30 and 2.64,
respectively). This signifies that shared burial space is the strongest
predictor of biological kinship in this dataset. Regarding biological
predictors, the Male-Male (XY-XY) combination stands out with a
consistent estimate of 1.94, significantly higher than the Female-Male
(XX-XY) estimate of \(\approx 0.90\). Furthermore, the Both Subadult
category (1.51 and 1.46) suggests a strong horizontal kinship signal
among children.

The multilevel hyperparameters reveal substantial variation at the
individual level (\(sd \approx 2.29\)), which suggests that the density
of kinship ties is not uniform across the population, but rather
concentrated around specific individuals. At the tomb level, the
estimated standard deviation increases from 1.16 in the Separated Model
to 1.86 in the Grouped Model. This increase suggests that grouping the
Elateia (T46, T56, and T62) tombs allows the model to better account for
differences between burial sites. Crucially, the fixed effects remained
consistent across both models, proving that the results are not
dependent on how the tombs were categorized.

\subsection{Model Interpretation}\label{model-interpretation}

The model output is reported in log-odds, as is standard in logistic
regression. However, coefficients expressed on the log-odds scale are
not directly intuitive to interpret. To make the results more
accessible, we therefore calculate marginal predicted probabilities.
These convert the model estimates into probabilities of the outcome,
which in this case, the probability that a pair of individuals is
biologically related.

\textbf{Marginal Effect (Predicted Probabilities)}:

Marginal predicted probabilities indicate how each explanatory variable
affects the average probability of relatedness, while holding all other
variables constant (ceteris paribus). This approach allows for a more
accessible interpretation of the model results by expressing effects in
terms of probabilities rather than log-odds.

??? me plots

The average predicted probabilities in the figure highlight several key
patterns across sex pairings, age compositions, and burial contexts.
Female--female pairs show relatively low average predicted probabilities
of relatedness (0.21\% in the separate tombs model and 0.31\% in the
grouped tombs model, ceteris paribus), while male--male pairs exhibit
higher probabilities compared to other sex combinations. For age
composition, pairs consisting of two subadults display higher predicted
probabilities than mixed or adult--adult pairs. The most pronounced
effect, however, is observed for burial context: pairs interred in the
same tomb show substantially higher predicted probabilities of
relatedness than any other category in both model specifications

\subsection{Model Robustness Check}\label{model-robustness-check}

To ensure that the results obtained in the primary MM-GLMM are not
artifacts of specific data distributions, prior specifications, or
classification choices, we performed three sensitivity analyses. These
tests evaluate whether our findings, particularly the dominant role of
the tomb effect and sex-specific patterns, remain consistent under
alternative model conditions.

\subsubsection{Sensitivity to Prior
Specification}\label{sensitivity-to-prior-specification}

In the primary model, we utilized flat (uninformative) priors, which
allow the data to drive the estimates entirely. To ensure the model was
not ``noise chasing'', which means fitting patterns that do not truly
exist, we re-fitted the model using Weakly Informative Priors
\(N(0,1)\). This approach applies a form of regularization, pulling
estimates toward zero unless the data provides strong evidence to the
contrary.

As expected with Priors \(N(0,1)\), the coefficients experienced a
slight shrinkage toward zero. No coefficients changed direction, and the
variable Same Tomb Indicator remained the dominant predictor, largely
unaffected by the prior. This confirms that the observed effects are
driven by strong signals in the data rather than being an artifact of
the prior choice.

\subsubsection{Sensitivity to the ``Same Tomb'' Predictor (Within-Tomb
Analysis)}\label{sensitivity-to-the-same-tomb-predictor-within-tomb-analysis}

Because the Same Tomb Indicator emerged as an exceptionally dominant
predictor in our primary model, it was necessary to conduct a robustness
check to ensure that other biological and demographic variables were not
merely statistical artifacts of this effect.

To isolate the influence of sex and age from the overwhelming signal of
shared burial space, we fitted a secondary model restricted exclusively
to within-tomb pairs. This approach effectively removes the across-tomb
variance, providing a focused look at the internal organization of the
burial units. The model was restricted to the 2,990 observations
involving individuals interred in the same tomb, compared to the 13,690
observations in the full dataset. Consequently, the estimated effects
appear weaker, with coefficients shrinking toward zero.

Despite the reduced power, the qualitative results remained consistent.
The direction of all effects was unchanged, and the Male-Male (XY-XY)
pair combination remained the dominant predictor of relatedness within
the tomb. This confirms that the sex-biased patterns identified in the
primary analysis are a fundamental characteristic of the kinship
structure within these tombs and are not driven by the broader ``Same
Tomb'' effect.

\subsubsection{Redefining the Relatedness Threshold (3rd Degree as
Unrelated)}\label{redefining-the-relatedness-threshold-3rd-degree-as-unrelated}

In our primary analysis, we collapsed 1st, 2nd, and 3rd-degree relatives
into a single ``Related'' category. However, 3rd-degree relationships
represent more distant biological ties. We conducted a robustness check
by redefining the response variable \(Y_{ij}\) to count only 1st and
2nd-degree relatives as ``Related.''

Under this stricter definition, the ``Same Tomb'' effect showed a
substantially larger estimated coefficient. This suggests that shared
burial context is an even stronger predictor for immediate kinship than
for distant relatives. From an archaeological perspective, 3rd-degree
kinmay not have resided in the same immediate household or been interred
in the family crypt as strictly as parent-child pairs. By including
3rd-degree relatives in the primary ``Related'' category, we effectively
``diluted'' the interment effect.

\section{Discussion}\label{discussion}

\subsection{Model assumption}\label{model-assumption}

While the MM-GLMM provides a robust framework for analyzing pairwise
kinship, the validity of its inferences depends on the underlying
statistical assumptions. This section discusses the theoretical and
empirical justifications for our model choices, as well as the potential
impacts of deviations from these assumptions.

\subsubsection{Equal Weighting of Random Effects (w =
0.5)}\label{equal-weighting-of-random-effects-w-0.5}

In the multi-membership structure, we assumed equal weights
(\(w_i = w_j = 0.5\) and \(w_{T(i)} = w_{T(j)} = 0.5\)) for the random
effects of individuals and tombs, respectively. This assumption is
theoretically plausible at the individual level, as there is no a priori
reason to believe one individual in a pair contributes more to the
shared probability of relatedness than the other. Furthermore,
simulation studies by \citet{smith2014impact} suggest that
multi-membership models are generally robust to weight misspecification,
meaning that even if the ``true'' weights were slightly unequal, the
resulting fixed-effect estimates remain reliable.

\subsubsection{Normality of Random
Effects}\label{normality-of-random-effects}

A standard assumption of GLMMs is that the random effects (\(u\) and
\(v\)) are independently and identically distributed following a normal
distribution: \(u,v \sim \mathcal{N}(0, \sigma^2)\). We evaluated this
through Q-Q plots for both the Individual and Tomb random effects across
our model variants.

For individual random effect diagnostic plots, the distribution is
somewhat skewed. Non-normality in random effects can potentially lead to
a biased global intercept and an inaccurate estimation of the random
effect distribution's shape. However, literature on mixed models
\citet{mcculloch2011misspecifying} demonstrates that fixed-effect
estimates and variance components are stable even when the normality
assumption for random effects is violated.

The diagnostic for the tomb-level effects is limited by the low number
of levels. According to \citet{maas2005sufficient}, the estimation of
the random effect variance (\(\sigma^2_{tomb}\) can be biased or
imprecise when the number of clusters is small (typically \(<30\)).
However the fixed-effect regression coefficients remain unbiased. The
model therefore remains a valid tool to interpret the influence of the
fixed-effect variables.

??? Put QQ Plots

\subsection{Selection Bias}\label{selection-bias}

A threshold of 15,000 overlapping SNPs was applied across all kinship
degrees to address the issue of missing not at random (MNAR) data. By
restricting the analysis to observations that meet this quality
criterion, the focus is placed exclusively on a high-quality subset of
the available genetic data. Although this procedure increases the
reliability of kinship estimates, it may introduce selection bias, as
individuals with lower data quality are systematically excluded from the
analysis. As a result, the findings may not be fully generalizable to
the broader population under study.

\subsection{Outlook for Future
Studies}\label{outlook-for-future-studies}

Future research should aim to expand the empirical basis of the analysis
by incorporating a substantially larger number of tombs across multiple
archaeological contexts. A broader sample would allow for the estimation
of robust tomb-level or multilevel models with a sufficient number of
higher-level units. This would enable a more reliable assessment of
tomb-specific characteristics---such as length of use, number of
interred individuals, and skeletal preservation patterns---and their
influence on genetic relatedness structures. With increased statistical
power, contextual variation between burial units could be modeled more
systematically.

In addition, future research should further address the issue of missing
data. In the present study, applying a single SNP threshold helped
reduce bias in the classification of kinship degrees. However, this
decision also meant that only individuals with sufficiently high genetic
data quality were included in the analysis. As a result, part of the
available information had to be excluded. Future studies could benefit
from improved sequencing quality and broader genomic coverage, which
would reduce the need for strict thresholds and allow more individuals
to be retained in the analysis. Moreover, applying alternative
statistical approaches to evaluate how different missing data
assumptions affect the results could help assess the stability and
robustness of the findings.

From a methodological perspective, future research could explore
extensions of the multi-membership generalized linear mixed model.
Relaxing assumptions such as equal random-effect weights or normally
distributed random effects may allow for a more flexible representation
of individual- and tomb-level heterogeneity. Furthermore, larger
datasets would permit more stable estimation of variance components,
particularly at the tomb level, thereby improving inference regarding
contextual effects.

Finally, future work should evaluate the potential impact of selection
processes introduced by quality thresholds. Comparing results obtained
from high-quality subsamples with those from broader inclusion criteria
or alternative imputation strategies may help assess the
generalizability of the findings to the wider population under study.

\section{A Acknowledgment}\label{a-acknowledgment}

We would like to express our sincere gratitude to Irene Högner and
Alissa Mittnik for their valuable collaboration and support as our
project partners. Their expertise, insightful feedback, and continuous
engagement significantly contributed to the development and successful
completion of this project.

We are especially grateful to our supervisor, Yichen Han, for the
dedicated guidance, constructive comments, and academic support
throughout the entire project. His supervision greatly helped us refine
our methodology, strengthen our analysis, and improve the overall
quality of this work.

\section{B Use of AI tools}\label{b-use-of-ai-tools}

\#\#\#Report ChatGPT (GPT-5.2) were used to support the drafting and
revision of the written report. The tools assisted in improving clarity,
structure, and linguistic precision, particularly in reformulating
paragraphs, refining academic wording, and ensuring coherence between
sections. All conceptual content, interpretations, statistical
reasoning, and conclusions were developed independently by the authors.
The AI tool did not generate original scientific arguments but served as
a language and structure support tool. All outputs were critically
reviewed, edited, and verified by the authors. ??Latex??

\#\#\#Code ChatGPT (GPT-5.2) and Gemini? were used to assist in
troubleshooting, refining, and optimizing statistical code (e.g., for
logistic regression models and data visualization). This included
suggestions for improving code structure, debugging errors, and
clarifying model implementation. However, all analytical decisions,
model specifications, variable selections, and interpretations of
results were made independently. The correctness and functionality of
all code were manually tested and validated before inclusion in the
analysis.

\#\#\#Literature Research ChatGPT (GPT-5.2) was used to support the
development of literature search strategies and to clarify theoretical
concepts. All academic sources cited in the report were independently
identified, accessed through scholarly databases, and critically
evaluated. No references were included without independent verification.
The interpretation and integration of the literature were conducted
manually.

\bibliography{ajs-template.bib}



\end{document}
